\section{Two notions of visibility}
\label{sec:visibility}

The examples separate two senses in which a ``nilpotent direction'' can be
visible, and keeping them apart explains exactly which algebras the two theories
admit.

\begin{definition}\label{def:visibility}
Let $a\in A$.
\begin{itemize}
\item $a$ is \emph{Jacobson-invisible} if $a\in\rad(A)$, i.e.\ $a$ is killed by
  every irreducible representation.  This is the scheme-theoretic sense of an
  infinitesimal thickening, as in $\RR[\varepsilon]/(\varepsilon^2)$.
\item $a$ is \emph{commutator-invisible} if $\gauge_1(a)=0$, i.e.\
  $a\in\Zc(A)$.
\end{itemize}
\end{definition}

The two agree on central elements (both invisible) and on matrix nilpotents in a
simple block (both visible), but they diverge in general.

\begin{proposition}[$T_2(\RR)$ separates the notions]\label{prop:visibility-T2}
In $T_2(\RR)$ the element $E_{12}$ lies in the radical, hence is
Jacobson-invisible --- it is killed by both diagonal characters --- yet
$\gauge_1(E_{12})=4>0$, so it is commutator-visible.  Consequently
$\gauge_1$ tracks non-centrality, which coincides with representational
visibility only when $\rad(A)\subseteq\Zc(A)$.
\end{proposition}

\begin{remark}[A claim that needs the constant]\label{rem:excess}
The earlier notes summarise the situation as: the excess of $\PB$ over the
intersection $\mathcal{I}$ is exactly the non-central radical, because the
$C^*$-instrument $a\mapsto a^*a$ forbids radical outright while $\gauge_1$
forbids only \emph{central} radical.  The first half of that is right --- central
radical is unfunded, hence always excluded.  The second half is not a theorem,
and at $C=\tfrac14$ it is contradicted by the one case we can compute:
$T_2(\RR)$ has non-central radical and is \emph{not} an object at $\tfrac14$
(Proposition~\ref{prop:T2-constant}).  What is true is that $\gauge_1$ forbids
central radical at every constant, and forbids the non-central radical of
$T_2(\RR)$ below $\kap^*$.  Whether it forbids \emph{all} radical at
$C=\tfrac14$ is Conjecture~\ref{conj:semisimple}.
\end{remark}

\begin{proposition}[$\gauge_\infty$ discriminates but is fragile]
\label{prop:dinf-fragile}
The spectral gauge excludes $T_2(\RR)$: every commutator there is a multiple of
the nilpotent $E_{12}$, so $\gauge_\infty\equiv0$ on $T_2(\RR)$ while the square
defect of $E_{12}$ is $1$, which is therefore unfunded.  It admits $M_2(\RR)$,
where $\gauge_\infty(E_{12})=1$ via the non-nilpotent commutator
$[E_{12},E_{21}]=E_{11}-E_{22}$.  This is a genuinely spectral,
semisimplicity-sensitive cut, matching the $C^*$ exclusion of the radical.  But
$\gauge_\infty$ vanishes on whole radical-commutator directions, and quotients
can manufacture such directions, so a $\gauge_\infty$-class has no reason to be
closed under quotients.
\end{proposition}

\begin{remark}[The trade-off, tabulated]
\begin{center}
\begin{tabular}{lll}
\toprule
Gauge & Binding budget on $T_2(\RR)$ & Behaviour \\
\midrule
$\gauge_1$ & $\gauge_1(E_{12})=4$, defect funded (but $\kap=\kap^*>\tfrac14$) & closure-friendly; base is functorial \\
large finite $q$ & $\gauge_q$ small but nonzero & plausibly closed; $T_2$ approximately out \\
$\gauge_\infty$ & $\gauge_\infty\equiv0$, defect unfunded & $T_2$ excluded; closure fails \\
\bottomrule
\end{tabular}
\end{center}
The first row's parenthetical is where the earlier notes and the present
treatment part company: at $C=\tfrac14$ the norm gauge \emph{does} exclude
$T_2(\RR)$ (Theorem~\ref{thm:constant-separates}), so the first row buys
discrimination as well as closure, and the trade-off the table describes may be
an artefact of looking only at constants above $\kap^*$.  The middle row is the
hope that a single large finite $q$ gets the discrimination of $\gauge_\infty$
without its zeros.  It is
Problem~\ref{prob:finite-q}, and we stress that
Lemma~\ref{lem:gauge-basic}(e) makes it \emph{not} merely a matter of tuning:
already at $q=2$ the vanishing locus is strictly larger than the centre, so the
functoriality of \S\ref{sec:base} is lost at every finite $q\ge2$, not only in
the limit.  Any positive answer to Problem~\ref{prob:finite-q} therefore buys
discrimination at the price of the geometry.
\end{remark}
